\section{Complete discrete valuation rings and lifting idempotents}

    Let \(R\) be a complete discrete valuation ring with normalized valuation
    \[
        v:K^\times \longrightarrow \mathbb{Z},
    \]
    where \(K=\operatorname{Frac}(R)\). Let \(\pi\in R\) be a uniformizer, so that
    \[
        v(\pi)=1.
    \]
    Denote the maximal ideal of \(R\) by
    \[
        P=J(R)=(\pi),
    \]
    and denote the residue field by
    \[
        F=R/P.
    \]

\begin{lemma}
    We have
    \[
        \bigcap_{n\geq 1} P^n=0.
    \]
\end{lemma}

\begin{proof}
    Since \(R\) is a discrete valuation ring, every nonzero element \(x\in R\) can be
    written uniquely in the form
    \[
        x=\pi^m u,
    \]
    where \(m\geq 0\) and \(u\in R^\times\). Hence \(x\in P^m\), but
    \[
        x\notin P^{m+1}.
    \]
    Therefore no nonzero element can lie in every \(P^n\). Thus
    \[
        \bigcap_{n\geq 1}P^n=0.
    \]
\end{proof}

\begin{proposition}
    Let \(V\) be a finitely generated \(R\)-module. Then \(V\) is a finite direct sum
    of cyclic \(R\)-modules. Thus we may write
    \[
        V=Ru_1\oplus \cdots \oplus Ru_r.
    \]
    For every \(n\geq 0\), one has
    \[
        P^nV=\pi^n V
        =
        (\pi^n R)u_1\oplus \cdots \oplus (\pi^n R)u_r.
    \]
    Moreover,
    \[
        \bigcap_{n\geq 1}P^nV=0.
    \]
\end{proposition}

\begin{proof}
    Since \(R\) is a discrete valuation ring, it is a principal ideal domain. Hence, by
    the structure theorem for finitely generated modules over a principal ideal domain,
    \(V\) is a finite direct sum of cyclic \(R\)-modules:
    \[
        V=Ru_1\oplus \cdots \oplus Ru_r.
    \]
    Therefore, for every \(n\geq 0\),
    \[
        P^nV=(\pi^n)V
        =
        (\pi^n R)u_1\oplus \cdots \oplus (\pi^n R)u_r.
    \]

    Since
    \[
        \bigcap_{n\geq 1}(\pi^n)=0,
    \]
    it follows componentwise that
    \[
        \bigcap_{n\geq 1}P^nV=0.
    \]
\end{proof}

\begin{definition}[\(P\)-adic metric]
    Let \(V\) be a finitely generated \(R\)-module. For \(0\neq x\in V\), define
    \[
        \nu_V(x):=\max\{\,n\geq 0\mid x\in P^nV\,\}.
    \]
    This is well-defined because
    \[
        \bigcap_{n\geq 1}P^nV=0.
    \]
    We also set
    \[
        \nu_V(0)=+\infty.
    \]
    Fix a real number \(0<c<1\). The \(P\)-adic metric on \(V\) is defined by
    \[
        d_P(x,y)=c^{\nu_V(x-y)}.
    \]
\end{definition}

\begin{theorem}
    Let \(V\) be a finitely generated \(R\)-module. If \(R\) is complete with respect
    to the \(P\)-adic metric, then \(V\) is complete with respect to the \(P\)-adic metric.
\end{theorem}

\begin{proof}
    Since \(V\) is finitely generated over the discrete valuation ring \(R\), it is a
    finite direct sum of cyclic \(R\)-modules:
    \[
        V=Ru_1\oplus \cdots \oplus Ru_r.
    \]
    Each cyclic summand \(Ru_i\) is isomorphic either to \(R\) or to \(R/(\pi^{m_i})\)
    for some \(m_i\geq 1\). Since \(R\) is complete, and quotients of complete
    \(P\)-adic modules by closed submodules are complete, each cyclic summand \(Ru_i\)
    is complete.

    A finite direct sum of complete metric \(R\)-modules is complete. Hence \(V\) is
    complete with respect to the \(P\)-adic metric.
\end{proof}

    Now let \(A\) be an \(R\)-algebra which is finitely generated as an \(R\)-module.
    For \(n\geq 0\), define
    \[
        A_n:=P^nA=\pi^n A.
    \]
    Since \(A\) is a finitely generated \(R\)-module and \(R\) is complete, the
    \(R\)-module \(A\) is complete with respect to the \(P\)-adic metric.

\begin{proposition}
    The multiplication map
    \[
        A\times A\longrightarrow A,
        \qquad
        (a,b)\longmapsto ab
    \]
    is continuous with respect to the \(P\)-adic topology.
\end{proposition}

\begin{proof}
    It is enough to check continuity at an arbitrary point \((a,b)\in A\times A\).
    Let \(\ell\geq 0\). Suppose
    \[
        a'-a\in P^\ell A,
        \qquad
        b'-b\in P^\ell A.
    \]
    Then
    \[
        a'b'-ab
        =
        (a'-a)b+a(b'-b)+(a'-a)(b'-b).
    \]
    Since \(P^\ell A\) is a two-sided ideal of \(A\), we have
    \[
        (a'-a)b\in P^\ell A,
        \qquad
        a(b'-b)\in P^\ell A,
    \]
    and also
    \[
        (a'-a)(b'-b)\in P^{2\ell}A\subseteq P^\ell A.
    \]
    Hence
    \[
        a'b'-ab\in P^\ell A.
    \]
    Therefore multiplication is continuous.
\end{proof}

\begin{proposition}
    The quotient
    \[
        \overline{A}:=A/A_1=A/PA
    \]
    is a finite-dimensional \(F\)-algebra, where
    \[
        F=R/P.
    \]
\end{proposition}

\begin{proof}
    Since \(A\) is finitely generated as an \(R\)-module, choose generators
    \[
        a_1,\dots,a_m\in A.
    \]
    Then their images
    \[
        \overline{a}_1,\dots,\overline{a}_m\in A/PA
    \]
    generate \(A/PA\) as an \(F=R/P\)-vector space. Hence \(A/PA\) is finite-dimensional
    over \(F\).

    The multiplication on \(A\) induces a multiplication on \(A/PA\), because \(PA\) is
    a two-sided ideal of \(A\). Therefore
    \[
        A/PA
    \]
    is a finite-dimensional \(F\)-algebra.
\end{proof}

\begin{theorem}
    Let \(A\) be a finitely generated \(R\)-algebra. Then:
    \begin{enumerate}
        \item
        \[
            A\pi \subseteq J(A).
        \]
        \item There exists an integer \(n>0\) such that
        \[
            J(A)^n \subseteq A\pi.
        \]
    \end{enumerate}
\end{theorem}

\begin{proof}
    Let \(V\) be a simple left \(A\)-module. Since \(A\pi\) is a two-sided ideal of
    \(A\), the subspace
    \[
        \pi V
    \]
    is an \(A\)-submodule of \(V\). Hence, by simplicity,
    \[
        \pi V=V
        \qquad\text{or}\qquad
        \pi V=0.
    \]

    Choose \(0\neq v\in V\). Since \(V\) is simple, we have
    \[
        V=Av.
    \]
    Since \(A\) is finitely generated over \(R\), it follows that \(V\) is finitely
    generated as an \(R\)-module.

    By Nakayama's lemma, because
    \[
        \pi\in J(R),
    \]
    the equality \(\pi V=V\) is impossible. Therefore
    \[
        \pi V=0.
    \]
    Thus \(A\pi\) annihilates every simple \(A\)-module. Hence
    \[
        A\pi \subseteq J(A).
    \]

    This proves \((1)\).

    For \((2)\), set
    \[
        A^\ast := A/A\pi.
    \]
    Since \(A\) is finitely generated as an \(R\)-module, \(A^\ast\) is a finite-dimensional
    algebra over the residue field
    \[
        F=R/(\pi).
    \]
    Hence \(A^\ast\) is an Artinian ring. Therefore its Jacobson radical is nilpotent.

    Since \(A\pi\subseteq J(A)\), we have
    \[
        J(A^\ast)=J(A/A\pi)=J(A)/A\pi.
    \]
    Hence there exists some \(n>0\) such that
    \[
        \bigl(J(A)/A\pi\bigr)^n=0.
    \]
    Equivalently,
    \[
        J(A)^n\subseteq A\pi.
    \]
    This proves \((2)\).
\end{proof}

\begin{theorem}
    Let \(R\) be a complete discrete valuation ring with uniformizer \(\pi\), and let
    \(A\) be an \(R\)-algebra which is finitely generated as an \(R\)-module. Let
    \(I\) be a two-sided ideal of \(A\) such that
    \[
        I\subseteq J(A).
    \]
    Write
    \[
        \overline{A}:=A/I.
    \]
    Let \(\overline{c}\in \overline{A}\) be an idempotent, and suppose
    \[
        \overline{c}
        =
        \overline{c}_1+\cdots+\overline{c}_n
    \]
    is an idempotent decomposition in \(\overline{A}\). Then:
    \begin{enumerate}
        \item The idempotent \(\overline{c}\) lifts to an idempotent \(e\in A\), and
        there exist orthogonal idempotents
        \[
            e_1,\dots,e_n\in A
        \]
        such that
        \[
            \overline{e_i}=\overline{c_i},
            \qquad 1\leq i\leq n,
        \]
        and
        \[
            e=e_1+\cdots+e_n.
        \]

        \item An idempotent \(e\in A\) is primitive if and only if its image
        \[
            \overline{e}\in \overline{A}
        \]
        is primitive.

        \item If \(e,f\in A\) are primitive idempotents, then
        \[
            Ae\cong Af
            \qquad\Longleftrightarrow\qquad
            \overline{A}\,\overline{e}\cong \overline{A}\,\overline{f}.
        \]
    \end{enumerate}
\end{theorem}

\begin{proof}
    Since \(I\subseteq J(A)\), by the previous theorem there exists an integer
    \(m>0\) such that
    \[
        I^m\subseteq A\pi.
    \]
    Hence, for every \(r\geq 1\),
    \[
        I^{mr}\subseteq A\pi^r.
    \]
    For each \(r\geq 1\), set
    \[
        A^{(r)}:=A/I^{mr}.
    \]

    We first prove that every idempotent of \(\overline{A}=A/I\) lifts to an
    idempotent of \(A\). Since \(I/I^m\) is a nilpotent ideal of \(A/I^m\), the
    idempotent \(\overline{c}\in A/I\) lifts to an idempotent in \(A/I^m\). Thus
    there exists \(c^{(1)}\in A\) such that
    \[
        \bigl(c^{(1)}\bigr)^2\equiv c^{(1)} \pmod{I^m},
        \qquad
        c^{(1)}\equiv c \pmod I.
    \]

    Next, since \(I^m/I^{2m}\) is a nilpotent ideal of \(A/I^{2m}\), the idempotent
    \(c^{(1)}+I^m\in A/I^m\) lifts to an idempotent in \(A/I^{2m}\). Hence there
    exists \(c^{(2)}\in A\) such that
    \[
        \bigl(c^{(2)}\bigr)^2\equiv c^{(2)} \pmod{I^{2m}},
        \qquad
        c^{(2)}\equiv c^{(1)} \pmod{I^m}.
    \]

    Repeating this argument, we obtain a sequence \(\{c^{(r)}\}_{r\geq 1}\) in \(A\)
    such that
    \[
        \bigl(c^{(r)}\bigr)^2\equiv c^{(r)} \pmod{I^{rm}},
    \]
    and
    \[
        c^{(r)}\equiv c^{(r-1)} \pmod{I^{(r-1)m}}
        \qquad (r\geq 2).
    \]
    Since
    \[
        I^{(r-1)m}\subseteq A\pi^{r-1},
    \]
    the sequence \(\{c^{(r)}\}_{r\geq 1}\) is Cauchy with respect to the
    \(\pi\)-adic topology on \(A\).

    As \(A\) is finitely generated over the complete ring \(R\), the \(R\)-module
    \(A\) is complete. Therefore the limit
    \[
        e:=\lim_{r\to\infty} c^{(r)}
    \]
    exists in \(A\). Since multiplication
    \[
        A\times A\longrightarrow A
    \]
    is continuous, passing to the limit in
    \[
        \bigl(c^{(r)}\bigr)^2-c^{(r)}\longrightarrow 0
    \]
    gives
    \[
        e^2=e.
    \]
    Moreover, \(e\equiv c\pmod I\), so \(e\) lifts \(\overline{c}\).

    We now lift the idempotent decomposition. Let \(e\in A\) be any idempotent lifting
    \(\overline{c}\). We construct, for each \(r\geq 1\), elements
    \[
        c_1^{(r)},\dots,c_n^{(r)}\in A
    \]
    such that their images in \(A^{(r)}=A/I^{rm}\) are orthogonal idempotents satisfying
    \[
        e\equiv c_1^{(r)}+\cdots+c_n^{(r)}
        \pmod{I^{rm}},
    \]
    and
    \[
        c_i^{(r)}\equiv c_i^{(r-1)}
        \pmod{I^{(r-1)m}}
        \qquad (r\geq 2).
    \]

    This is obtained inductively by applying the lifting theorem for idempotent
    decompositions modulo nilpotent ideals. Indeed, the kernel of
    \[
        A/I^{rm}\longrightarrow A/I^{(r-1)m}
    \]
    is
    \[
        I^{(r-1)m}/I^{rm},
    \]
    which is nilpotent.

    For each fixed \(i\), the sequence
    \[
        \{c_i^{(r)}\}_{r\geq 1}
    \]
    is Cauchy, because
    \[
        c_i^{(r)}-c_i^{(r-1)}\in I^{(r-1)m}\subseteq A\pi^{r-1}.
    \]
    Hence the limit
    \[
        e_i:=\lim_{r\to\infty}c_i^{(r)}
    \]
    exists in \(A\).

    Passing to the limit, and using the continuity of multiplication, we get
    \[
        e_i^2=e_i,
        \qquad
        e_ie_j=0\quad (i\neq j),
    \]
    and
    \[
        e=e_1+\cdots+e_n.
    \]
    Moreover,
    \[
        \overline{e_i}=\overline{c_i}.
    \]
    This proves \((1)\).

    We prove \((2)\). Recall that \(J(A)\) contains no nonzero idempotents. Indeed, if
    \(x\in J(A)\) and \(x^2=x\), then \(1-x\) is a unit, and
    \[
        x(1-x)=0
    \]
    implies \(x=0\).

    Suppose first that \(\overline{e}\) is primitive. If
    \[
        e=e_1+e_2
    \]
    is an idempotent decomposition in \(A\), then reducing modulo \(I\) gives
    \[
        \overline{e}
        =
        \overline{e_1}+\overline{e_2}.
    \]
    Since \(\overline{e}\) is primitive, one of \(\overline{e_1}\), \(\overline{e_2}\)
    is zero. Suppose \(\overline{e_1}=0\). Then
    \[
        e_1\in I\subseteq J(A).
    \]
    Since \(e_1\) is an idempotent and \(J(A)\) contains no nonzero idempotents, we
    get
    \[
        e_1=0.
    \]
    Hence \(e\) is primitive.

    Conversely, suppose that \(e\) is primitive. If \(\overline{e}\) were not primitive,
    then there would be a nontrivial idempotent decomposition
    \[
        \overline{e}
        =
        \overline{d}_1+\overline{d}_2
    \]
    in \(\overline{A}\). By part \((1)\), this decomposition lifts to an idempotent
    decomposition
    \[
        e=d_1+d_2
    \]
    in \(A\). Since \(\overline{d}_1\neq 0\) and \(\overline{d}_2\neq 0\), both \(d_1\)
    and \(d_2\) are nonzero. This contradicts the primitiveness of \(e\). Therefore
    \(\overline{e}\) is primitive.

    This proves \((2)\).

    Finally, we prove \((3)\). Since \(I\subseteq J(A)\), for every idempotent \(e\in A\)
    we have
    \[
        Ie\subseteq J(A)e=\operatorname{Rad}(Ae).
    \]
    Therefore the natural epimorphism
    \[
        Ae\longrightarrow Ae/Ie
    \]
    is essential. Moreover,
    \[
        Ae/Ie\cong \overline{A}\,\overline{e}.
    \]
    Since \(Ae\) is projective, the map
    \[
        Ae\longrightarrow \overline{A}\,\overline{e}
    \]
    is a projective cover.

    Similarly,
    \[
        Af\longrightarrow \overline{A}\,\overline{f}
    \]
    is a projective cover.

    If
    \[
        Ae\cong Af,
    \]
    then passing to the quotient modulo \(I\) gives
    \[
        \overline{A}\,\overline{e}
        \cong
        \overline{A}\,\overline{f}.
    \]

    Conversely, if
    \[
        \overline{A}\,\overline{e}
        \cong
        \overline{A}\,\overline{f},
    \]
    then \(Ae\) and \(Af\) are projective covers of isomorphic modules. By uniqueness
    of projective covers, we obtain
    \[
        Ae\cong Af.
    \]

    Hence
    \[
        Ae\cong Af
        \qquad\Longleftrightarrow\qquad
        \overline{A}\,\overline{e}\cong \overline{A}\,\overline{f}.
    \]
\end{proof}
